\pagestyle{empty}
\tight
\begin{tblr}{width=\textwidth,colspec={|Q[c,m,1.9cm]|X[l,m]|},hlines,vlines,hspan=minimal,row{1}={bg=headblue},rowsep=1pt,colsep=3pt}
\SetCell{c,m}\hfil\logo\hfil & \textbf{POLITEKNIK NEGERI MALANG}\par\textbf{JURUSAN TEKNOLOGI INFORMASI}\par\textbf{PROGRAM STUDI: D4 TEKNIK INFORMATIKA} \\
\SetCell[c=2]{c} \textbf{RENCANA PEMBELAJARAN SEMESTER (RPS)} & \\
\end{tblr}
\begin{longtblr}{width=\textwidth,colspec={|Q[l,m,1.9cm]|X[0.85,l,m]|X[1.45,l,m]|X[0.75,l,m]|X[1.15,l,m]|},hlines,vlines,hspan=minimal,rowsep=1pt,colsep=2pt,row{1,3}={bg=headgray,font=\bfseries}}
MATA KULIAH & KODE & BOBOT (SKS)/jam & SEMESTER & TGL. PENYUSUNAN \\
Kewarganegaraan & RTI253003 & 2 SKS/3 jam & 3 & 29 Januari 2026 \\
OTORISASI & Dosen Pengembang RPS & Koordinator RMK & \SetCell[c=2]{l} Ka PRODI &  \\
 & Dr. Widaningsih Condrowardani, S.Psi, SH, MH & Dr. Widaningsih Condrowardani, S.Psi, SH, MH & \SetCell[c=2]{l}{\vspace{8mm}\par Dr. Ely Setyo Astuti, S.T., M.T.} &  \\
\SetCell[r=11]{l} Capaian Pembelajaran (CP) & \SetCell[c=4]{l,bg=headgray,font=\bfseries} Capaian Pembelajaran Lulusan Program Studi (CPL-Prodi) &  &  &  \\
 & CPL01 & \SetCell[c=3]{l} Bertakwa kepada Tuhan Yang Maha Esa dan mampu menunjukkan sikap religius, menjunjung tinggi nilai kemanusiaan, hukum, disiplin, serta etika profesi; berpartisipasi dalam lifelong learning; dan peka terhadap isu sosial-teknologi. &  &  \\
 & \SetCell[c=4]{l,bg=headgray,font=\bfseries} Capaian Pembelajaran mata kuliah (CPL-MK) &  &  &  \\
 & CPMK0102 & \SetCell[c=3]{l} Mahasiswa mampu menginternalisasi nilai ketakwaan, Pancasila, kewarganegaraan, dan tanggung jawab sosial sebagai landasan sikap profesional serta berkomitmen pada lifelong learning dan adaptasi karir. &  &  \\
 & \SetCell[c=4]{l,bg=headgray,font=\bfseries} Kemampuan akhir tiap tahapan belajar (Sub-CPMK) &  &  &  \\
 & SCPMK0102-02101 & \SetCell[c=3]{l} Mahasiswa mampu menjelaskan konsep identitas nasional, negara, konstitusi, dan hubungan negara dengan warga negara sebagai fondasi kehidupan berbangsa dan bernegara. &  &  \\
 & SCPMK0102-02102 & \SetCell[c=3]{l} Mahasiswa mampu menganalisis sistem demokrasi dan negara hukum dalam konteks kehidupan berbangsa dan bernegara di Indonesia. &  &  \\
 & SCPMK0102-02103 & \SetCell[c=3]{l} Mahasiswa mampu menjelaskan konsep hak asasi manusia, prinsip-prinsipnya, dan mekanisme penegakannya dalam kehidupan berbangsa dan bernegara. &  &  \\
 & SCPMK0102-02104 & \SetCell[c=3]{l} Mahasiswa mampu menguraikan wawasan nusantara, ketahanan nasional, dan integrasi nasional sebagai implementasi nilai-nilai kewarganegaraan dalam konteks global. &  &  \\
\SetCell[r=2]{l} Penilaian & \SetCell[c=4]{l} *Nilai CPMK didapat dari agregasi nilai SUB-CPMK yang dibebankan &  &  &  \\
 & \SetCell[c=4]{l}{\tiny\begin{tblr}{width=\linewidth,colspec={|Q[c,m,0.1\linewidth]|Q[c,m,0.16\linewidth]|X[l,m]|X[l,m]|X[l,m]|X[l,m]|X[l,m]|Q[c,m,0.08\linewidth]|},hlines,vlines,hspan=minimal,rowsep=1pt,colsep=0.7pt,row{1,2}={bg=headgray,font=\bfseries}}
Kode CPMK & Kode SCPMK & \SetCell[c=5]{c} Bobot per Bentuk Penilaian & & & & & Total Bobot \\
 & & Tugas & Kuis & UTS & PBL & UAS & \\
\SetCell[r=1]{c} \SetCell[r=4]{c} CPMK0102 & SCPMK0102-02101 & 8\%: Tugas Mgg 1--6, 9--13 & 4\%: Kuis Mgg 7 & 10\%: UTS Mgg 8 & & 10\%: UAS Mgg 17 & 32\% \\
\SetCell[r=3]{c}  & SCPMK0102-02102 & 6\%: Tugas Mgg 6, 9--10 & 3\%: Kuis Mgg 7 & 8\%: UTS Mgg 8 & 5\%: PBL Mgg 14 & 10\%: UAS Mgg 17 & 32\% \\
 & SCPMK0102-02103 & 3\%: Tugas Mgg 11--12 & 2\%: Kuis Mgg 7 & 1\%: UTS Mgg 8 & 5\%: PBL Mgg 15 & 8\%: UAS Mgg 17 & 19\% \\
 & SCPMK0102-02104 & 3\%: Tugas Mgg 13 & 1\%: Kuis Mgg 7 & 1\%: UTS Mgg 8 & 5\%: PBL Mgg 16 & 7\%: UAS Mgg 17 & 17\% \\
\SetCell[c=7]{r}\textbf{Total Per Penilaian} & & & & & & & \textbf{100\%} \\
\end{tblr}} &  &  &  \\
Diskripsi Singkat Mata Kuliah & \SetCell[c=4]{l} Mata kuliah Kewarganegaraan membangun kesadaran kebangsaan dan cinta tanah air melalui kajian identitas nasional, konstitusi, demokrasi, hak asasi manusia, wawasan nusantara, ketahanan nasional, dan integrasi nasional. Mahasiswa mengembangkan sikap demokratis, disiplin, dan partisipatif berdasarkan sistem nilai Pancasila. &  &  &  \\
Bahan Kajian / Materi Pembelajaran & \SetCell[c=4]{l}\begin{enumerate}\item Identitas Nasional\item Negara dan Konstitusi\item Hubungan Negara dan Warga Negara\item Negara Hukum\item Demokrasi\item Hak Asasi Manusia\item Wawasan Nusantara\item Ketahanan Nasional\item Integrasi Nasional\end{enumerate} &  &  &  \\
\SetCell[r=2]{l} Pustaka & \SetCell[c=4]{l} \textbf{Utama:}\begin{enumerate}\item Widaningsih. 2023. Pendidikan Kewarganegaraan. Yogyakarta: Pustaka Egaliter.\end{enumerate} &  &  &  \\
 & \SetCell[c=4]{l} \textbf{Pendukung:}\begin{enumerate}\item UPT MKU-Politeknik Negeri Malang. 2016. Pendidikan Kewarganegaraan. Malang: Aditya Media Publishing.\item Azra, A. 2003. Demokrasi Hak Asasi Manusia dan Masyarakat Madani. Jakarta: Prenata Media.\item Russel, B. 2004. Sejarah Filsafat Barat dan Kaitannya dengan Kondisi Sosio-Politik dari Zaman Kuno hingga Sekarang. Yogyakarta: Pustaka Pelajar.\end{enumerate} &  &  &  \\
Media Pembelajaran & \SetCell[c=2]{l} \textbf{Software:} Microsoft Word, Microsoft PowerPoint, LMS. &  & \SetCell[c=2]{l} \textbf{Hardware:} Komputer, laptop, proyektor. &  \\
Nama Dosen Pengampu & \SetCell[c=4]{l}\begin{enumerate}\item Dr. Widaningsih Condrowardani, S.Psi, SH, MH\end{enumerate} &  &  &  \\
Matakuliah Syarat & \SetCell[c=4]{l} - &  &  &  \\
\end{longtblr}
\begin{longtblr}{width=\textwidth,colspec={|Q[c,m,0.06\textwidth]|X[1.35,l,m]|X[1.08,l,m]|X[1.16,l,m]|X[0.7,l,m]|X[1.24,l,m]|X[1.06,l,m]|X[0.98,l,m]|Q[c,m,0.05\textwidth]|},hlines,vlines,hspan=minimal,rowhead=2,row{1,2}={bg=headgreen,font=\bfseries},rowsep=1pt,colsep=1.5pt}
Mgg.\par Ke & Kemampuan Akhir Yang Direncanakan (Sub-CP-MK) & Bahan kajian (Materi Pembelajaran) & Bentuk dan Metode Pembelajaran & Estimasi Waktu & Pengalaman Belajar Mahasiswa & Kriteria \& Bentuk Penilaian & Indikator Penilaian & Bobot Penilaian (\%) \\
(1) & (2) & (3) & (4) & (5) & (6) & (7) & (8) & (9) \\
1 & SCPMK0102-02101 & Identitas nasional: pengertian identitas nasional dan identitas nasional sebagai karakter bangsa. & Modalitas: Blended Learning. Bentuk: Luring/Daring. Metode: Ceramah, diskusi kelompok, studi kasus, dan tanya jawab. & 1 x 3 x 50' tatap muka; tugas/belajar mandiri. & Mahasiswa menyimak, berdiskusi, dan menganalisis pengertian identitas nasional serta mengerjakan tugas analisis kewarganegaraan. & Kriteria: kejelasan dan ketepatan analisis. Bentuk: tugas tertulis analisis kewarganegaraan. Mengacu rubrik RTM. & Ketepatan pemahaman konsep; kualitas analisis tertulis; partisipasi diskusi. & 1 \\
2 & SCPMK0102-02101 & Identitas nasional: proses berbangsa dan bernegara, serta politik identitas. & Modalitas: Blended Learning. Bentuk: Luring/Daring. Metode: Ceramah, diskusi kelompok, studi kasus, dan tanya jawab. & 1 x 3 x 50' tatap muka; tugas/belajar mandiri. & Mahasiswa menyimak, berdiskusi, dan menganalisis proses berbangsa dan bernegara serta mengerjakan tugas analisis kewarganegaraan. & Kriteria: kejelasan dan ketepatan analisis. Bentuk: tugas tertulis analisis kewarganegaraan. Mengacu rubrik RTM. & Ketepatan pemahaman konsep; kualitas analisis tertulis; partisipasi diskusi. & 2 \\
3 & SCPMK0102-02101 & Negara dan konstitusi: pengertian negara, pengertian konstitusi, dan peranan konstitusi dalam kehidupan bernegara. & Modalitas: Blended Learning. Bentuk: Luring/Daring. Metode: Ceramah, diskusi kelompok, studi kasus, dan tanya jawab. & 1 x 3 x 50' tatap muka; tugas/belajar mandiri. & Mahasiswa menyimak, berdiskusi, dan menganalisis negara dan konstitusi serta mengerjakan tugas analisis kewarganegaraan. & Kriteria: kejelasan dan ketepatan analisis. Bentuk: tugas tertulis analisis kewarganegaraan. Mengacu rubrik RTM. & Ketepatan pemahaman konsep; kualitas analisis tertulis; partisipasi diskusi. & 2 \\
4 & SCPMK0102-02101 & Hubungan negara dan warga negara 1: pengertian hak dan kewajiban, serta hak dan kewajiban warga negara menurut UUD 1945. & Modalitas: Blended Learning. Bentuk: Luring/Daring. Metode: Ceramah, diskusi kelompok, studi kasus, dan tanya jawab. & 1 x 3 x 50' tatap muka; tugas/belajar mandiri. & Mahasiswa menyimak, berdiskusi, dan menganalisis hak dan kewajiban warga negara serta mengerjakan tugas analisis kewarganegaraan. & Kriteria: kejelasan dan ketepatan analisis. Bentuk: tugas tertulis analisis kewarganegaraan. Mengacu rubrik RTM. & Ketepatan pemahaman konsep; kualitas analisis tertulis; partisipasi diskusi. & 2 \\
5 & SCPMK0102-02101 & Hubungan negara dan warga negara 2: pelaksanaan hak dan kewajiban, serta esensi dan urgensi harmoni hak dan kewajiban negara dan warga negara. & Modalitas: Blended Learning. Bentuk: Luring/Daring. Metode: Ceramah, diskusi kelompok, studi kasus, dan tanya jawab. & 1 x 3 x 50' tatap muka; tugas/belajar mandiri. & Mahasiswa menyimak, berdiskusi, dan menganalisis harmoni hak dan kewajiban serta mengerjakan tugas analisis kewarganegaraan. & Kriteria: kejelasan dan ketepatan analisis. Bentuk: tugas tertulis analisis kewarganegaraan. Mengacu rubrik RTM. & Ketepatan pemahaman konsep; kualitas analisis tertulis; partisipasi diskusi. & 2 \\
6 & SCPMK0102-02102 & Negara hukum: pengertian dan ciri negara hukum, serta makna Indonesia sebagai negara hukum. & Modalitas: Blended Learning. Bentuk: Luring/Daring. Metode: Ceramah, diskusi kelompok, studi kasus, dan tanya jawab. & 1 x 3 x 50' tatap muka; tugas/belajar mandiri. & Mahasiswa menyimak, berdiskusi, dan menganalisis konsep negara hukum serta mengerjakan tugas analisis kewarganegaraan. & Kriteria: kejelasan dan ketepatan analisis. Bentuk: tugas tertulis analisis kewarganegaraan. Mengacu rubrik RTM. & Ketepatan pemahaman konsep; kualitas analisis tertulis; partisipasi diskusi. & 2 \\
7 & SCPMK0102-02101, SCPMK0102-02102 & Kuis Konsep Kewarganegaraan: materi minggu 1--6 (identitas nasional, negara dan konstitusi, hubungan negara dan warga negara, negara hukum). & Modalitas: Luring. Bentuk: Kuis tertulis. Metode: Ujian individu. & 1 x 3 x 50' tatap muka. & Mahasiswa mengerjakan soal kuis konsep kewarganegaraan secara mandiri. & Kriteria: ketepatan jawaban. Bentuk: kuis tertulis. & Ketepatan jawaban sesuai kunci. & 10 \\
8 & SCPMK0102-02101, SCPMK0102-02102 & UTS: materi minggu 1--7 (identitas nasional, negara dan konstitusi, hubungan negara dan warga negara, negara hukum). & Modalitas: Luring. Bentuk: Ujian tertulis. Metode: Ujian individu, esai dan/atau pilihan ganda. & 1 x 3 x 50'. & Mahasiswa menjawab soal UTS materi minggu 1--7 secara mandiri dan jujur. & Kriteria: ketepatan dan kedalaman jawaban. Bentuk: ujian tertulis esai/pilihan ganda. & Ketepatan jawaban; kejelasan argumentasi. & 20 \\
9 & SCPMK0102-02102 & Demokrasi 1: konsep dasar demokrasi, serta prinsip-prinsip dan indikator demokrasi. & Modalitas: Blended Learning. Bentuk: Luring/Daring. Metode: Ceramah, diskusi kelompok, studi kasus, dan tanya jawab. & 1 x 3 x 50' tatap muka; tugas/belajar mandiri. & Mahasiswa menyimak, berdiskusi, dan menganalisis konsep dasar demokrasi serta mengerjakan tugas analisis kewarganegaraan. & Kriteria: kejelasan dan ketepatan analisis. Bentuk: tugas tertulis analisis kewarganegaraan. Mengacu rubrik RTM. & Ketepatan pemahaman konsep; kualitas analisis tertulis; partisipasi diskusi. & 2 \\
10 & SCPMK0102-02102 & Demokrasi 2: perjalanan demokrasi di Indonesia, pendidikan demokrasi, dan pentingnya demokrasi sebagai sistem politik kenegaraan modern. & Modalitas: Blended Learning. Bentuk: Luring/Daring. Metode: Ceramah, diskusi kelompok, studi kasus, dan tanya jawab. & 1 x 3 x 50' tatap muka; tugas/belajar mandiri. & Mahasiswa menyimak, berdiskusi, dan menganalisis perjalanan demokrasi di Indonesia serta mengerjakan tugas analisis kewarganegaraan. & Kriteria: kejelasan dan ketepatan analisis. Bentuk: tugas tertulis analisis kewarganegaraan. Mengacu rubrik RTM. & Ketepatan pemahaman konsep; kualitas analisis tertulis; partisipasi diskusi. & 2 \\
11 & SCPMK0102-02103 & Hak Asasi Manusia 1: pengertian HAM, prinsip-prinsip HAM, dan sejarah perkembangan HAM. & Modalitas: Blended Learning. Bentuk: Luring/Daring. Metode: Ceramah, diskusi kelompok, studi kasus, dan tanya jawab. & 1 x 3 x 50' tatap muka; tugas/belajar mandiri. & Mahasiswa menyimak, berdiskusi, dan menganalisis konsep dasar HAM serta mengerjakan tugas analisis kewarganegaraan. & Kriteria: kejelasan dan ketepatan analisis. Bentuk: tugas tertulis analisis kewarganegaraan. Mengacu rubrik RTM. & Ketepatan pemahaman konsep; kualitas analisis tertulis; partisipasi diskusi. & 2 \\
12 & SCPMK0102-02103 & Hak Asasi Manusia 2: jenis-jenis HAM, dasar hukum HAM nasional dan internasional, serta mekanisme penegakan dan pelanggaran HAM. & Modalitas: Blended Learning. Bentuk: Luring/Daring. Metode: Ceramah, diskusi kelompok, studi kasus, dan tanya jawab. & 1 x 3 x 50' tatap muka; tugas/belajar mandiri. & Mahasiswa menyimak, berdiskusi, dan menganalisis mekanisme penegakan HAM serta mengerjakan tugas analisis kewarganegaraan. & Kriteria: kejelasan dan ketepatan analisis. Bentuk: tugas tertulis analisis kewarganegaraan. Mengacu rubrik RTM. & Ketepatan pemahaman konsep; kualitas analisis tertulis; partisipasi diskusi. & 2 \\
13 & SCPMK0102-02104 & Wawasan Nusantara 1: wilayah sebagai ruang hidup dan konsep wawasan nusantara sebagai penerapan geopolitik Indonesia. & Modalitas: Blended Learning. Bentuk: Luring/Daring. Metode: Ceramah, diskusi kelompok, studi kasus, dan tanya jawab. & 1 x 3 x 50' tatap muka; tugas/belajar mandiri. & Mahasiswa menyimak, berdiskusi, dan menganalisis konsep wawasan nusantara serta mengerjakan tugas analisis kewarganegaraan. & Kriteria: kejelasan dan ketepatan analisis. Bentuk: tugas tertulis analisis kewarganegaraan. Mengacu rubrik RTM. & Ketepatan pemahaman konsep; kualitas analisis tertulis; partisipasi diskusi. & 1 \\
14 & SCPMK0102-02104 & Wawasan Nusantara 2: unsur-unsur dasar wawasan nusantara, penerapan wawasan nusantara, dan tantangan implementasinya. & Modalitas: Blended Learning. Bentuk: Luring/Daring. Metode: Project Based Learning (PBL), diskusi kelompok, presentasi. & 1 x 3 x 50' tatap muka; tugas/belajar mandiri. & Mahasiswa mengerjakan proyek berbasis masalah terkait implementasi wawasan nusantara secara berkelompok, berdiskusi, dan mempresentasikan hasil. & Kriteria: kedalaman analisis dan kualitas proyek. Bentuk: proyek kelompok dan presentasi. Mengacu rubrik RTM. & Ketepatan analisis masalah; kualitas solusi proyek; kualitas presentasi. & 5 \\
15 & SCPMK0102-02104 & Ketahanan Nasional: pengertian dan sejarah ketahanan nasional, unsur-unsur ketahanan nasional, pendekatan asta gatra, serta globalisasi dan ketahanan nasional. & Modalitas: Blended Learning. Bentuk: Luring/Daring. Metode: Project Based Learning (PBL), diskusi kelompok, presentasi. & 1 x 3 x 50' tatap muka; tugas/belajar mandiri. & Mahasiswa mengerjakan proyek berbasis masalah terkait ketahanan nasional secara berkelompok, berdiskusi, dan mempresentasikan hasil. & Kriteria: kedalaman analisis dan kualitas proyek. Bentuk: proyek kelompok dan presentasi. Mengacu rubrik RTM. & Ketepatan analisis masalah; kualitas solusi proyek; kualitas presentasi. & 5 \\
16 & SCPMK0102-02104 & Integrasi Nasional: integrasi nasional dan pluralitas masyarakat Indonesia, dinamika integrasi nasional, strategi integrasi, serta integrasi nasional Indonesia. & Modalitas: Blended Learning. Bentuk: Luring/Daring. Metode: Project Based Learning (PBL), diskusi kelompok, presentasi. & 1 x 3 x 50' tatap muka; tugas/belajar mandiri. & Mahasiswa mengerjakan proyek berbasis masalah terkait integrasi nasional secara berkelompok, berdiskusi, dan mempresentasikan hasil. & Kriteria: kedalaman analisis dan kualitas proyek. Bentuk: proyek kelompok dan presentasi. Mengacu rubrik RTM. & Ketepatan analisis masalah; kualitas solusi proyek; kualitas presentasi. & 5 \\
17 & SCPMK0102-02101, SCPMK0102-02102, SCPMK0102-02103, SCPMK0102-02104 & UAS: komprehensif seluruh materi minggu 1--16 (identitas nasional, negara dan konstitusi, hubungan negara dan warga negara, negara hukum, demokrasi, HAM, wawasan nusantara, ketahanan nasional, integrasi nasional). & Modalitas: Luring. Bentuk: Ujian tertulis komprehensif. Metode: Ujian individu, esai dan/atau pilihan ganda. & 1 x 3 x 50'. & Mahasiswa menjawab soal UAS komprehensif materi minggu 1--16 secara mandiri dan jujur. & Kriteria: ketepatan dan kedalaman jawaban. Bentuk: ujian tertulis esai/pilihan ganda. & Ketepatan jawaban; kejelasan argumentasi; pemahaman komprehensif. & 35 \\
\end{longtblr}
